\documentclass[11pt,a4paper]{article}
\usepackage[margin=22mm]{geometry}
\usepackage{amsmath,amssymb,booktabs,microtype}
\usepackage[hidelinks]{hyperref}
\newcommand{\ii}{\mathrm i}
\newcommand{\ket}[1]{|#1\rangle}
\newcommand{\bra}[1]{\langle#1|}
\newcommand{\spec}{\operatorname{spec}}
\title{Route G-R9b: Paired Exchange Channels, an Exact $4/9$ Bound,\\and Two Ways to One Outlet}
\author{A side branch of G-R9, alongside G-R10; no mass, time-law or flux fit}
\date{9 October 2026}
\begin{document}
\maketitle
\begin{abstract}
G-R9 found that a fixed vertex contact between two identical three-node
rings opens a second, degenerate exchange outlet. Here we show that this
is structural. Flipping the sign of every link maps $L$ to $4-L$; on a
three-node ring this pairs the lower gap of one ring with the upper gap of
its mirror partner. Hence \emph{every} resonant pair of flux 3-rings has
the reverse outlet degenerate, and every real density--density contact
couples the two outlets with equal strength. For the rank-one vertex
contact, two exact dark states make the single-outlet probability at most
$4/9$ at every update and every contact strength, which G-R9 had stated
only as a weak-contact envelope. One outlet is reached in two ways. A
3-ring paired with a 4-ring, joined by the same vertex contact, leaves only
the target degenerate and transfers with probability $1-O(\gamma^2)$. Two
identical rings joined by an \emph{oriented} contact form an exactly closed
chiral triangle that transfers with unit probability. In both, the two
copies must be distinguishable. Ordered updates, quantum mechanics, the
resonant flux and the contact orientation remain inputs.
\end{abstract}

\section{What G-R9 left open}
G-R9 used $L_\theta=2I-e^{\ii\theta/3}S-e^{-\ii\theta/3}S^\dagger$ on two
copies, read clock A on modes $(0,1)$ and clock B on modes $(1,2)$, and
asked for the exchange $\ket{1,1}\to\ket{0,2}$. At $\theta=\pi/2$ the free
gaps are both $\sqrt3$. The contact $\gamma P_0\otimes P_0$ then also
reaches $\ket{2,0}$, which has the same energy. Swap symmetry gave an exact
bound $P_{11\to02}\le\tfrac12$; a weak-contact block gave the envelope
$4/9$, which was explicitly \emph{not} asserted for the full generator.
Its next bounded choice was either to keep all three channels or to find
an independent reason that removes the extra outlet. G-R10 took the first
branch: with the contact unchanged it built the all-channel readout,
found two exact dark states and a conserved $2/3$ dark weight of
$\ket{1,1}$. This side note, written in parallel and not replacing G-R10,
asks whether the extra outlet is an accident of identical copies at
$\theta=\pi/2$, and gives one construction for one-outlet exchange in each
direction of that choice.

\paragraph{Notation.} On an $N$-node ring, $S\ket j=\ket{j+1\bmod N}$ and
\begin{equation}
 L_N(\theta)=2I-e^{\ii\theta/N}S-e^{-\ii\theta/N}S^\dagger,\qquad
 \lambda^{(N)}_k(\theta)=2-2\cos\frac{\theta-2\pi k}{N},
\end{equation}
with flux-independent Fourier eigenvectors $f_k(j)=e^{2\pi\ii jk/N}/\sqrt N$.
Labels $a,b$ below are energy-sorted modes of each copy. For a 3-ring with
$0<\theta<\pi$ the sorted order is $k=0,1,2$. The input is $\ket{1,1}$, the
\emph{target} outlet $\ket{0,2}$ and the \emph{reverse} outlet $\ket{2,0}$.

\section{A link sign flip pairs the channels}
\paragraph{Lemma 1 (link sign flip).}
$4I-L_N(\theta)=L_N(\theta+N\pi)$; $L_N(\theta+2\pi)=DL_N(\theta)D^\dagger$
with $D={\rm diag}(e^{2\pi\ii j/N})$; and $L_N(\theta)^*=L_N(-\theta)$.
\emph{Proof.} Replacing every link amplitude by its negative adds $\pi$ to
each of the $N$ link phases, so the loop phase grows by $N\pi$. The other
two identities follow from $DSD^\dagger=e^{2\pi\ii/N}S$ and $S^*=S$.
\hfill$\square$

For even $N$ the shift $N\pi$ is a multiple of $2\pi$, so the spectrum is
symmetric about 2 at every flux (the ring is bipartite). For $N=3$ the
shift is $\pi$ modulo $2\pi$, and with conjugation
\begin{equation}
 \spec L_3(\pi-\theta)=4-\spec L_3(\theta).\label{mirror}
\end{equation}
The sorted gaps follow in closed form:
\begin{equation}
 \Delta_-(\theta)=\lambda_1-\lambda_0=2\sqrt3\,\sin\frac{\pi-\theta}{3},\qquad
 \Delta_+(\theta)=\lambda_2-\lambda_1=2\sqrt3\,\sin\frac{\theta}{3}.\label{gaps}
\end{equation}

\paragraph{Theorem 1 (channel pairing).}
Let A and B be 3-rings with fluxes $\theta_A,\theta_B\in(0,\pi)$. Then
$E_{02}=E_{11}$ holds if and only if $\theta_B=\pi-\theta_A$, and in that
case also $E_{20}=E_{11}$.
\emph{Proof.} $E_{02}=E_{11}$ means $\Delta_-(\theta_A)=\Delta_+(\theta_B)$.
By Eq.~\eqref{gaps} this is $\sin\frac{\pi-\theta_A}{3}=\sin\frac{\theta_B}{3}$
with both arguments in $(0,\pi/3)$, so $\theta_B=\pi-\theta_A$. Then
$E_{20}-E_{11}=\Delta_+(\theta_A)-\Delta_-(\pi-\theta_A)=0$.\hfill$\square$

So no choice of the two fluxes separates the outlets. G-R9's identical
copies at $\pi/2$ are the self-mirror case of this family. Any flux is
gauge- or conjugation-equivalent to one in $[0,\pi]$, and at the ends a
ring level is itself degenerate. The pairing is special to three levels:
for $N>3$ the flip pairs gap $(0,1)$ with gap $(N-2,N-1)$, not with the
middle dial. The verifier finds the resonant partner by bisection on 49
values of $\theta_A$ and confirms $\theta_B=\pi-\theta_A$ and
$E_{20}=E_{11}$ to $10^{-10}$.

\section{What a local contact can and cannot see}
Take a real density--density contact
$V=\sum_{j,l}c_{jl}\,P_j\otimes P_l$ between two 3-rings. In the mode basis
\begin{equation}
 \bra{a,b}V\ket{a',b'}=\hat c(a'-a,\,b'-b),\qquad
 \hat c(p,q)=\frac19\sum_{j,l}c_{jl}\,\omega^{jp+lq},\quad \omega=e^{2\pi\ii/3}.
\end{equation}
\paragraph{Lemma 2 (equal coupling).}
On the resonant triple $(\ket{0,2},\ket{1,1},\ket{2,0})$ every diagonal
entry is $c_0=\hat c(0,0)$, and
$V_{02,11}=V_{11,20}=V_{20,02}=z\equiv\hat c(1,-1)$. The block is
therefore circulant, $c_0I+zC+z^*C^\dagger$: a three-site ring whose loop
phase is $\Phi=3\arg z$. Both outlets couple to the input with the same
magnitude $|z|$. If $c$ is symmetric, $z$ is real and $\Phi\in\{0,\pi\}$.
\emph{Proof.} The three elements $V_{02,11}$, $V_{11,20}$ and $V_{20,02}$
all carry momentum transfer $(1,-1)$ modulo 3. Reality of $c$ gives
$\hat c(-p,-q)=\hat c(p,q)^*$, so $V_{20,11}=z^*$.\hfill$\square$

With real $z$ the block is $c_0-z+3z\ket T\bra T$, with
$\ket T=(\ket{0,2}+\ket{1,1}+\ket{2,0})/\sqrt3$. In the weak-contact block this is G-R9's $4/9$
structure, now for every symmetric contact and every mirror pair. For
identical rings a symmetric $c$ also makes the full generator
swap-symmetric, so G-R9's exact bound $P\le\frac12$ covers all of them.

\paragraph{Lemma 3 (rank-one dark states).}
Let $H=H_0+\gamma\ket v\bra v$ and let $\mathcal D_0$ be the eigenspace of
$H_0$ at $E_0$, containing $\ket i\neq\ket t$. Put
$\ket b=P_{\mathcal D_0}\ket v/\|P_{\mathcal D_0}v\|$,
$w_x=|\langle x|b\rangle|^2$ and $G(s)=\bra b e^{-\ii Hs}\ket b$. Then for
every $s$
\begin{equation}
 \bra t e^{-\ii Hs}\ket i=\langle t|b\rangle\langle b|i\rangle
 \big[G(s)-e^{-\ii E_0s}\big],\qquad P_{i\to t}(s)\le4w_iw_t.\label{dark}
\end{equation}
\emph{Proof.} Vectors of $\mathcal D_0$ orthogonal to $\ket b$ are also
orthogonal to $\ket v$, so they are exact eigenvectors at $E_0$; their
orthogonal complement is invariant and contains $\ket b$. Write
$\ket i=\langle b|i\rangle\ket b+\ket{d_i}$ and use
$\langle t|d_i\rangle=-\langle t|b\rangle\langle b|i\rangle$. The bound
follows from $|G|\le1$.\hfill$\square$

The vertex contact is rank one with all $|v_{ab}|^2=1/9$. For identical
rings or any mirror pair, $w=1/3$ on the triple, so
\begin{equation}
 P_{11\to02}(s)=P_{11\to20}(s)=\tfrac19\big|G(s)-e^{-\ii E_0s}\big|^2\le\tfrac49
\end{equation}
\emph{exactly}, for every $\gamma$ and every update. The two outlets have
identical amplitudes even for mirror pairs, which lack swap symmetry.
Here $\ket{1,1}$ has bright weight $1/3$, the complement of G-R10's
conserved dark weight $2/3$. For weak contact
$G\simeq e^{-\ii(E_0+\gamma/3)s}$, which reproduces G-R9's
$\frac49\sin^2(\gamma s/6)$ and G-R10's weak-contact amplitude; Eq.~\eqref{dark}
makes that amplitude exact once the exact $G$ is used. The verifier checks Eq.~\eqref{dark} against
full $9\times9$ evolution, and computes $G$ independently from the roots of
the secular equation $1=\gamma\sum_{ab}|v_{ab}|^2/(E-\varepsilon_{ab})$.

\begin{center}
\begin{tabular}{rrrr}\toprule
$\theta_A$ & $\gamma$ & max sampled $P_{02}$ & bound\\\midrule
$\pi/2$ & 0.01 & 0.444444088 & 4/9\\
$\pi/2$ & 0.5 & 0.439608352 & 4/9\\
$\pi/2$ & 10 & 0.417446514 & 4/9\\
$\pi/3$ (B at $2\pi/3$) & 2 & 0.386793114 & 4/9\\\bottomrule
\end{tabular}
\end{center}
Equation~\eqref{dark} also says what unit transfer through a rank-one
contact requires: $w_i=w_t=\tfrac12$. The input's energy shell must contain
exactly one other state that the contact touches.

\section{Exit A: distinguishable spectra}
Keep the G-R9 vertex contact and replace copy B by a 4-ring. For
$0<\theta<\pi$ its sorted spectrum is
$2\mp2\cos(\theta/4)$, $2\mp2\sin(\theta/4)$, so its middle gap is
$4\sin(\theta/4)$. With A at $\theta_A=\pi/2$ ($\Delta_-=\sqrt3$),
resonance fixes
\begin{equation}
 \theta_B=4\arcsin\frac{\sqrt3}{4}=1.791329588,\qquad
 \spec L_4(\theta_B)=\Big\{2\mp\frac{\sqrt{13}}2,\;2\mp\frac{\sqrt3}2\Big\}.
\end{equation}
Among the twelve product levels, the shell of $\ket{1,1}$
($E_0=4-\sqrt3/2$) contains only $\ket{1,1}$ and $\ket{0,2}$. The nearest
other level is the reverse outlet, detuned by $(3\sqrt3-\sqrt{13})/2=0.7953$.

This holds across the whole resonant domain. A level $\ket{a,b}$ is
degenerate with $\ket{1,1}$ iff $\lambda_a-\lambda_1=\mu_1-\mu_b$. Since B
is non-degenerate, $a=1$ forces $b=1$ and $a=0$ forces $b=2$; only $a=2$,
$b=0$ could add a state. On a 2001-point grid of
$\theta_A\in(0.2756,\pi)$, where the lower end is the flux at which
$\Delta_-^A$ equals the largest possible 4-ring middle gap $2\sqrt2$, the
reverse detuning never falls below $0.318$.

Now Lemma~3 has $w_i=w_t=\tfrac12$ and no bound below one:
$P_{11\to02}=\tfrac14|G-e^{-\ii E_0s}|^2$. The bright pair
$(\ket{1,1}+\ket{0,2})/\sqrt2$ is shifted by $\gamma/6$ while
$(\ket{1,1}-\ket{0,2})/\sqrt2$ is an exact eigenstate, so for weak contact
$P\simeq\sin^2(\gamma s/12)$. With $\alpha=1/4$ and integer updates, the
full $12\times12$ evolution gives:
\begin{center}
\begin{tabular}{rrrrr}\toprule
$\gamma$ & best $n$ & $P_{02}$ & $P_{20}$ & $(1-P_{02})/\gamma^2$\\\midrule
0.09 & 856 & 0.999543095 & $9.7\times10^{-5}$ & 0.056\\
0.03 & 2535 & 0.999936858 & $4.3\times10^{-7}$ & 0.070\\
0.01 & 7558 & 0.999994927 & $3.1\times10^{-7}$ & 0.051\\\bottomrule
\end{tabular}
\end{center}
The best update is later than the weak estimate $6\pi/(\gamma\alpha)$
by about 2\%, 0.9\% and 0.2\%: the exact bright root is
$E_0+\tfrac\gamma6(1-0.196\gamma+\dots)$, so the exchange slows slightly. The deficit is $O(\gamma^2)$ and comes from dressing by
off-shell levels. The reverse outlet is not removed by selecting an
operator: the partner's spectrum detunes it, and the contact is unchanged.
What is still supplied is one tuned flux, as at G-R9's $\theta=\pi/2$.

\section{Exit B: an oriented contact on identical rings}
Keep two identical rings and orient the contact:
\begin{equation}
 V_h=\gamma\sum_{j,l}h\big((j-l)\bmod3\big)\,P_j\otimes P_l .
\end{equation}
It commutes with the joint translation $S\otimes S$, which conserves
$a+b\bmod3$. The class $a+b\equiv2$ is exactly the resonant triple, so the
triple is a closed sector for every $\gamma$. Unlike the vertex contact,
there is no leakage. Here $c_0=\gamma(h_0+h_1+h_2)/3$ and
$z=\tfrac\gamma3(h_0+h_1\omega+h_2\omega^2)$. Swapping the copies sends
$h(d)\to h(-d)$, so a swap-symmetric profile ($h_1=h_2$) has real $z$,
$\Phi=0$, and stays at $4/9$.

The profile $h=(2,0,1)$ links node $j$ of A to node $j$ of B with weight
$2\gamma$ and to node $j+1$ with weight $\gamma$. It gives
$z=(\gamma/\sqrt3)e^{-\ii\pi/6}$ and $\Phi=-\pi/2$. The circulant then has
equally spaced eigenvalues $c_0+\{\sqrt3,0,-\sqrt3\}|z|$, and
\begin{equation}
 e^{-\ii Ht_*}\ket{1,1}=e^{\ii\varphi}\ket{0,2},\qquad t_*=\frac{2\pi}{3\gamma},
\end{equation}
for some phase $\varphi$. This is the chiral quantum walk on a triangle: a loop phase of
$\pm\pi/2$ makes transport directional~\cite{zimboras}.
The reversed orientation $h=(2,1,0)$ sends the quantum to $\ket{2,0}$
instead. The same holds for any mirror pair. Sampled at integer updates
($\alpha=1/4$):
\begin{center}
\begin{tabular}{rrrrr}\toprule
$\gamma$ & $t_*$ & $n$ & $P_{02}$ at $n$ & $P_{20}$ at $n$\\\midrule
0.09 & 23.271 & 93 & 0.9999976057 & $1.2\times10^{-6}$\\
0.03 & 69.813 & 279 & 0.9999976057 & $1.2\times10^{-6}$\\
0.01 & 209.440 & 838 & 0.9999997561 & $1.2\times10^{-7}$\\\bottomrule
\end{tabular}
\end{center}
The deficit is only the offset between $t_*$ and the nearest update; at
$t_*$ itself the full $9\times9$ value is one to $10^{-10}$. Here the
third channel is not removed. The oriented contact routes the quantum
around the resonant triangle, and all three channels stay in the model.
What is supplied is an orientation and the weight ratio that sets
$\Phi=\pm\pi/2$. That value gives exact transfer at $t_*$; it is not
claimed to be the only value that can give high transfer at later times.

Exit B does \emph{not} remove G-R10's preparation obstruction. Any
translation-covariant contact makes the closed triple circulant, so
$\ket T$ is an exact eigenvector and the $\ket T$ weight $1/3$ of
$\ket{1,1}$ is conserved. G-R10's dark clock state
$\ket{D_+}=(\ket{0,2}-2\ket{1,1}+\ket{2,0})/\sqrt6$ is orthogonal to
$\ket T$, so its fidelity from $\ket{1,1}$ never exceeds $2/3$, for every
profile $h$ and every time. The verifier checks this for the oriented,
reversed, symmetric and four random profiles.

\section{What changed, and the next bounded step}
\begin{center}
\begin{tabular}{lccl}\toprule
copies and contact & states at $E_{11}$ & best single outlet & status\\\midrule
identical, vertex (G-R9) & 3 & $\le4/9$ & exact, all $\gamma$\\
mirror pair, vertex & 3 & $\le4/9$ & exact, all $\gamma$\\
identical, any symmetric $c$ & 3 & $\le1/2$ & exact (swap)\\
3-ring + 4-ring, vertex (A) & 2 & $1-O(\gamma^2)$ & full $12\times12$\\
identical, oriented (B) & 3 & $1$ at $t_*$ & exact closed sector\\\bottomrule
\end{tabular}
\end{center}

\paragraph{Physical reading.}
A one-for-one exchange between two dials needs the copies to be
distinguishable, either in their spectra (Exit A) or in how their contact
is oriented (Exit B). Identical 3-rings, or any mirror pair, joined by an
unoriented local contact split the quantum between two outlets: exactly
for identical copies and for the vertex contact, and in the weak-contact
block otherwise. In the
language of Chen's Sec.~1.9~\cite{chen}, exchange between compatible event
rhythms becomes one-to-one only when the partners differ in some
respect. This is our inference for this toy model, not a statement of the
source.

\paragraph{Still assumed.}
Quantum kinematics, the update order $n$, the common $\alpha$, the
preparation and the readout are inputs, as in G-R9. Exit A tunes one flux;
Exit B chooses an orientation and a ratio. No dynamical principle selects
either. No physical seconds, records, metric or particle spectrum follow.

\paragraph{Relation to G-R10 and next step.}
The two notes agree: the dark states behind the exact $4/9$ bound are the
same exact dark states that G-R10 uses as a stationary conditional clock.
Exits A and B are declared alternatives for one-outlet \emph{exchange}
only. Neither prepares G-R10's $\ket{D_+}$, so G-R10's own next choice
(a declared preparation and record instrument, or the unfiltered
stationary mixture) is unchanged. Exit B is an exactly closed, solvable
three-channel exchange on which G-R10's readout could be repeated.
Separately, a spatial carrier would need a gapless network rather than
another internal ring; that is not attempted here.

\paragraph{Executable verification.}
The companion program is
\texttt{verify\_gr9b\_paired\_channels.py}.
Its inputs are fixed in \texttt{PAIRED\_CHANNEL\_CARD.json}, and it uses
numpy only. Its 586 checks
cover: Lemma~1 as matrix identities for $N=3,4$; the mirror~\eqref{mirror}
and gaps~\eqref{gaps}; Theorem~1 by independent bisection; Lemma~2 on 12
random real contacts; Lemma~3 against full evolution and against the
secular-equation roots, with the $4/9$ bound at strong contact; the Exit-A
census, domain scan, second-order bright root and $12\times12$ fixtures;
and the Exit-B closure,
transfer, orientation reversal and the $2/3$ bound on $\ket{D_+}$. Propagators use eigendecomposition and
are cross-checked by an independent Taylor scaling-and-squaring
exponential. The G-R9 code, card and TeX are preserved by SHA-256. These
are checks of finite algebra, not experiments.

\begin{thebibliography}{9}
\bibitem{chen} BoYu Chen, \emph{Energy spectrum for elementary particles},
final version, 23 April 2022, section 1.9. Conceptual source only.
\bibitem{zimboras} Z. Zimbor\'as, M. Faccin, Z. K\'ad\'ar, J. D. Whitfield,
B. P. Lanyon and J. Biamonte, Quantum transport enhancement by
time-reversal symmetry breaking, Sci.\ Rep.\ 3, 2361 (2013);
\href{https://arxiv.org/abs/1208.4049}{arXiv:1208.4049}.
\end{thebibliography}
\end{document}
